\section{Limitations and threats to validity}
\label{sec:threats}

\paragraph{Statistical power.} Seventeen folds give limited power. The 95\,\% confidence interval
on the central contrast (pretrained versus random-init) spans roughly $\pm 3$ pp, so an advantage
smaller than that cannot be resolved. Our result is a failure to establish an advantage, not
evidence of its absence. Increasing folds requires more participants, not more seeds.

\paragraph{One dataset, one task, one model.} We audit MtReNIRS on a three-class auditory
block-design dataset. We did not evaluate GeoNIRS, nor a visual or affective paradigm. The
auditory paradigm may be unusually favourable to systemic signal, which would make the
short-separation result paradigm-specific. Generalisation across datasets and tasks is the
obvious next experiment and we do not claim it here.

\paragraph{The random-encoder control is a lower bound.} Our control re-initialises each module
with its own default initialisation and runs ten seeds. A stronger control would tune the random
encoder's readout as carefully as the pretrained one, which could raise it further; a weaker
control (single seed) could understate it. We report the seed distribution so the reader can
judge.

\paragraph{Probe hyperparameters are fixed.} All arms use $C=1$, PCA-64 and one solver setting.
We did not perform nested hyperparameter selection, because doing so identically for twelve arms
was not affordable in our compute budget. A tuned probe could change the ranking between arms
whose accuracies differ by a few points. Total dimensionality is matched, which removes the most
obvious artefact but not this one.

\paragraph{Preprocessing choices are defaults.} The modified Beer--Lambert constants, the
differential pathlength factor of 6.0, the 0.01--0.2~Hz band and the run-mean intensity reference
are defensible defaults, not optimised values. Because epochs are z-scored per participant and
channel, the short-separation ablation acts on waveform shape rather than absolute amplitude;
a pipeline that preserved amplitude might show a different pattern.

\paragraph{Contamination accounting is by published list.} We determine that the evaluation
dataset is absent from the pretraining corpus from the twelve dataset names published with the
model. If the dataset entered pretraining under an unrecognised name, the pretrained arm would
be advantaged, so this uncertainty works in the model's favour, not against it.

\paragraph{Inference-only audit.} We evaluate frozen embeddings and a linear probe; we do not
fine-tune the encoder. Fine-tuning could change the picture substantially and is a natural
follow-up. We also do not claim any result about the model's published benchmark numbers, which
we did not attempt to reproduce exactly, because our preprocessing and probe differ.

\paragraph{Compute and numerical detail.} All inference and probes ran on a laptop CPU. One
earlier run of the stacking ablation used a probe that produced convergence warnings; we re-ran
it under the final protocol and report only the re-run (the qualitative conclusion was
unchanged). One earlier run of the short-separation variant silently failed to change the data
because the short channels had been removed before regression; the current pipeline asserts that
the cleaning step changes the input, and we verified the effect numerically (correlation
$0.434 \to 0.000$). We mention both because they are ordinary failure modes of this kind of
pipeline and are worth guarding against.
